\section{Supplemental Results}

\subsection{Equivalence Ratio Dependence}

On 2023-05-24, we performed measurements with a varying equivalence ratio. 

\begin{itemize}
    \item 516 - a position sweep at a fixed equivalence ratio of 0.6
    \item 5x3\_pos - a position and equivalence ratio sweep at five equivalence ratios and 0.1\% potassium mass fraction. The lower potassium mass fraction allowed for longer times on each test case allowed for both AES $n_{K}$ and MWT measurements. 
    \item 5x6\_pos - a position and equivalence ratio sweep at five equivalence ratios and 1\% potassium mass fraction. The high potassium mass uses up seed quickly, leaving enough time for only MWT measurements. 
\end{itemize}
 
\subsubsection{5x3 and 5x6 Position Dependence Test Case Sequences}

This section contains additional measurements for the 5x3 and 5x6 test case sequences. These were two-dimensional parameter sweeps of both stage position and equivalence ratio. See Figure \ref{fig:SI_expt_overview_2023-05-24} for how the sequence was performed. 

Figure \ref{fig:SI_5x3_pos_nK_m3} shows the AES measurements for the 5x3\_pos test case. Note the calibration of the AES signal for this test case exhibited significant calibration drift (see Figure \ref{fig:SI_5x3_pos_alpha}), and systematic error may be present in the $nK$ values. 

%TODO: Make sure every mention of nk is formatted as $nk$

\begin{figure}[]
\centering
\includegraphics[width=0.8\textwidth]{\repodir/experiment/analysis/tcs/output/figures/5x3_pos_nK_m3.png}
\caption{AES measurements 5x3\_pos}
\label{fig:SI_5x3_pos_nK_m3}
\end{figure}

\begin{figure}[]
\centering
\includegraphics[width=0.8\textwidth]{\repodir/experiment/analysis/tcs/output/figures/5x3_pos_alpha.png}
\caption{Absorption data for AES measurements in 5x3\_pos. The absorption offset at large wavelengths is due to drift in calibration.}
\label{fig:SI_5x3_pos_alpha}
\end{figure}%mp = "motor position"?


% \begin{figure}[]
% \centering
% \includegraphics[width=0.8\textwidth]{\repodir/experiment/analysis/tcs/output/figures/5x3_pos_mwt_KOH.png}
% \caption{MWT measurements 5x3\_pos. CFD centerline values for KOH are also shown.}
% \label{fig:SI_5x3_pos_mwt_KOH}
% \end{figure}


% \begin{figure}[]
% \centering
% \includegraphics[width=0.8\textwidth]{\repodir/experiment/analysis/tcs/output/figures/5x6_pos_mwt_KOH.png}
% \caption{MWT measurements 5x6\_pos. CFD centerline values for KOH are also shown.}
% \label{fig:SI_5x6_pos_mwt_KOH}
% \end{figure}


\begin{figure}[]
\centering
\includegraphics[width=0.8\textwidth]{\repodir/experiment/analysis/tcs/output/figures/5x3_pos_AS.png}
\caption{Position-dependent pre-pulse (-50 to -1 us) AS average for 5x3\_pos, shown as: Top: Pre-pulse AS versus stage position, with color representing equivalence ratio, and Bottom: Pre-pulse AS represented as a color scale, versus both stage position and equivalence ratio.}
\label{fig:SI_5x3_pos_AS}
\end{figure}


\begin{figure}[]
    \centering
    \includegraphics[width=0.8\linewidth]{\repodir/experiment/analysis/tcs/output/figures/5x3_pos_dAS.png} 
    \caption{Position dependent $\Delta AS$ for 5x3\_pos, shown as: Top: $\Delta AS$ versus stage position, with color representing equivalence ratio, and Bottom: $\Delta AS$ represented as a color scale, versus both stage position and equivalence ratio.}
    \label{fig:SI_5x3_pos_dAS}
\end{figure}


\begin{figure}[]
\centering
\includegraphics[width=0.9\textwidth]{\repodir/experiment/analysis/tcs/output/figures/5x6_pos_mwt_AStime_phi.png}
\caption{$\Delta AS$ time profiles for 5x6\_pos, at two stage positions. Left: $\Delta AS$ time profile at stage position of 55.13 mm. Right: $\Delta AS$ time profile at stage position of 180.1 mm (SFR-maximized position).}
\label{fig:SI_5x6_pos_mwt_AStime_phi}
\end{figure}


\clearpage
\subsubsection{516 Position Dependence Test Case Sequence}

This section contains additional measurements for the 516 test case. This test case was a position sweep at fixed nominal equivalence ratio 0.6. The higher seed rate and longer time on each test case allows for more accurate determination of AS exponential time constants than 5x3\_pos and 5x6\_pos test case sequences. The AES measurements associated with this test case are presented in the main text Figure 3C. 

Figure \ref{fig:SI_536_516_pos_mwt_fitting} shows position-dependent MWT measurements for the 516\_pos and 536\_pos test cases. Figure \ref{fig:SI_tau_cfd_pos_phi_compare} shows the position-dependent CFD-predicted electron lifetimes for K+ and O2\_A for both the 516 and 536 test cases. 

In Figure \ref{fig:SI_536_516_pos_mwt_fitting}B, position dependent SFR is shown. For phi = 0.59, the SFR is maximized at 155 mm downstream (compared to 180 mm for phi=0.79). We therefore compare lifetime measurements for 0.59 equvalence ratio with positions 25 mm upstream compared to measurements at 0.8 equivalence ratio. In Figure \ref{fig:SI_536_516_pos_mwt_fitting}B, we calculate the ratio of $\tau_{exp}$ at the respective SFR-maximized positions for the two equivalence ratios. In Figure \ref{fig:SI_tau_ratio_cfd_pos}, we calculate the ratio of $\tau_{pred}$ of the CFD data between the two equivalence ratios, but with the 0.6 equivalence ratio data shifted 25 mm downstream. The ratio of $\tau_{pred}$ is similar for both K+ and O2 and is consistent with the experimental measurements.

\begin{figure}[]
\centering
\includegraphics[width=0.8\textwidth]{\repodir/final/figures/SI/output/536_516_pos_mwt_fitting.png}
\caption{516 and 536 MWT AS exponential time constant comparison. A) raw data and exponential fits. The 0.6 equivalence ratio data is from 2024-05-24 and the 0.8 equivalence ratio data is from 2024-05-18 Run 1. The fit range is from 3 to 15 microseconds, which is not the same as in the main text but was chosen in this case to provide reasonable fits for all positions. The numbers on the right-hand side of each plot indicate the stage position in mm. B): Signal to fluctuation ratio (SFR, see Eq.\ 9 in the main text) versus stage position. C) $\tau_{exp}$ versus stage position. The ratio of $\tau_{exp}$ for the 516 and 536 test cases at their respective SFR-maximized positions is shown.} 
\label{fig:SI_536_516_pos_mwt_fitting}
\end{figure}

%DN: Throughout this section, you refer to $\Delta AS$ as MWT AS. Are they the same thing? %LA: AS and $\Delta AS$ are different but both MWT measurements. 

\begin{figure}[]
\centering
\includegraphics[]{\repodir/final/analysis/output/figures/tau_cfd_pos_phi_compare.png}
\caption{Comparison of position-dependent $\tau_{pred}$ for K+ and O2\_A for both the 516 and 536 test cases. The lifetimes are calculated from centerline CFD species densities (See also Figure \ref{fig:SI_cfd_species_tau_pos}).}
\label{fig:SI_tau_cfd_pos_phi_compare}
\end{figure}



\begin{figure}[]
\centering
\includegraphics[]{\repodir/final/analysis/output/figures/tau_ratio_cfd_pos.png}
\caption{Ratio of $\tau_{pred}$ for K+ and O2\_A between 0.6 and 0.8 equivalence ratio. For 0.6 equivalence ratio, the values are shifted downstream to match the SFR-maximized position of the 0.8 equivalence ratio. This allows for comparison to the experimental ratio calculated in Figure \ref{fig:SI_536_516_pos_mwt_fitting}.}
\label{fig:SI_tau_ratio_cfd_pos}
\end{figure}

\clearpage
\subsection{ICCD Camera Measurements}

This section contains additional measurements by the ICCD camera. 

The ICCD camera is calibrated using a projective transformation. The calibration image is shown in Figure \ref{fig:SI_iccd_calibration}. The calibration is performed by taking a picture of a ruler held at the exit of the HVOF barrel. The parameters for the projective transformation are determined by comparing the length of the ruler at the exit of the barrel and at 7 inches. From the transformed image, the pixel to mm conversion calibration is determined, and applied to calibrate the images. 

\begin{figure}[]
\centering
\includegraphics[width=0.8\textwidth]{\repodir/final/analysis/topcam/output/figures/calibration_image_projective.png}
\caption{Calibration of the ICCD images. A) Raw image used for calibration B) A projective transformation is applied to the image. C) Checking the calibration pipeline on the original image, with vertical lines shown at each inch.}
\label{fig:SI_iccd_calibration}
\end{figure}

\begin{figure}
    \centering
    \includegraphics[width=0.8\linewidth]{\repodir/final/figure_panels/output/figures/536_iccd_img_gatedelay.png} 
    \caption{ICCD camera images taken at 1\% potassium mass fraction in the Run 53x on 2023-05-24. During this test case the camera acquired 15 gate delay sweeps, which are averaged for these images. A) Images are taken at different gate delays. The laser is position at the nominal position. }
    \label{fig:SI_536_iccd}
\end{figure}


\begin{figure}
    \centering
    \includegraphics[]{\repodir/final/figure_panels/output/figures/536_iccd_laserspot_zoom.png} 
    \caption{Zoomed in view of the laser spot in the ICCD images shown in Figure \ref{fig:SI_536_iccd}. The laser is positioned at the nominal position and averaged over all gate delays.}
    \label{fig:SI_536_iccd_laserspot_zoom}
\end{figure}


\begin{figure}[]
\centering
\includegraphics[width=0.8\textwidth]{\repodir/final/analysis/topcam/output/figures/spe_cfd_comparison_2023-05-18.png}
\caption{Comparison of ICCD images and CFD predictions. This figure is similar to Figure 2 in the main text, but for the camera intensity, the data is averaged over all gate delays and time throughout a position sweep, which has the effect of showing all positions at once, as well as reducing the noise for the sans-laser image.}
\label{fig:SI_iccd_cfd_comparison}
\end{figure}


\begin{figure}[]
\centering
\includegraphics[]{\repodir/final/analysis/topcam/output/figures/spe_cfd_comparison_2023-05-18_KOH_overlap.png}
\caption{Overlap of CFD KOH and ICCD laser on and laser off difference image. The laser on and laser off difference image is overlaid on the CFD KOH image with a transparency of 0.5. For the camera intensity, the data is averaged over all gate delays and time throughout a position sweep, which has the effect of showing all positions at once.}
\label{fig:SI_iccd_cfd_comparison_KOH_overlap}
\end{figure}

\begin{figure}
    \centering
    \includegraphics[width=0.8\linewidth]{\repodir/final/analysis/topcam/output/figures/spe_cfd_comparison_2023-05-18_shockdiamonds.png} 
    \caption{Examining shock diamonds from ICCD Images and CFD. Top to bottom: A) Centerline of the Top camera laser off intensity and B) CFD pressure}
    \label{fig:SI_iccd_cfd_comparison_KOH}
\end{figure}

\subsection{Laser power dependence}

\begin{figure}
    \centering
    \includegraphics[]{\repodir/experiment/analysis/tcs/output/figures/fit_mwt_decay_power.png} 
    \caption{Exponential time constant of the MWT signal as a function of laser power. The time constant is calculated from the mono-molecular regime of the MWT signal.}
    \label{fig:SI_fit_mwt_decay_power}
\end{figure}

\clearpage

\subsection{Supplemental CFD results}

\begin{figure}[h]
    \includegraphics[]{\repodir/final/analysis/output/figures/cfd_species_tau_pos.png} 
    \caption{Top: Centerline species concentrations predicted by CFD for 1\% seed and 0.8 equivalence ratio. "Yeq" denotes thermal equilibrium species concentrations calculated through a Gibbs free energy minimization (using the Cantera Python package) at the CFD-predicted temperature.  Bottom: electron lifetimes for recombination with the given species density. The SFR-maximized position is indicated with a vertical line.}
    \label{fig:SI_cfd_species_tau_pos}
\end{figure}



\begin{figure}[h]
    \includegraphics[]{\repodir/final/analysis/output/figures/cfd_species_charged_sigma_pos.png} 
    \caption{Top: Centerline charged species concentrations for 1\% seed and 0.8 equivalence ratio. Bottom: centerline electron conductivity versus position along the length of the flame.}
    \label{fig:SI_cfd_species_charged_sigma_pos}
\end{figure}



\clearpage

We define an effective oxygen bimolecular recombination rate, $k_{r,b,eff}$, measured from the experiment. This requires that we assume O2 attachment is the dominant recombination pathway during the exponential regime of the MWT decay. We define $k_{r,b,eff}$ as 

\begin{equation}
    \label{eq:krbeff}
    k_{r,b,eff} = \frac{1}{\tau_{experiment} *  n_{O2}},
\end{equation}
where $\tau_{experiment}$ is the measured exponential time constant,  and $n_{O2}$ is the oxygen number density predicted from CFD. Additionally, if we assume that the collisional recombination partner of this effective reaction is the entire gas mixture, we can calculate the effective thermomolecular recombination coefficient as


\begin{equation}
    \label{eq:krteff}
    k_{r,t,eff} = \frac{1}{\tau_{experiment} * n_{all} *  n_{O2}},
\end{equation}
where $n_{all}$ is the total number density predicted from CFD. Figure \ref{fig:SI_krb_eff_53x} shows the calculation of $k_{r,t,eff}$ for 53x.

\begin{figure}[]
\centering
\includegraphics[width=0.7\textwidth]{\repodir/final/analysis/mwt/output/figures/krb_eff_53x.png}
\caption{Effective recombination coefficient calculation for the 53x test case sequence A) The exponential time constant measured from the MWT mono-molecular regime, averaged over all runs. B) The CFD-predicted total and O2 number density. C) The effective bimolecular recombination coefficient calculated from Equation \ref{eq:krbeff}} D) The effective termolecular rate calculation from Equation \ref{eq:krteff}
\label{fig:SI_krb_eff_53x}
\end{figure}

\subsection{MWT Supplemental Results}

\begin{figure}[h]
    \centering
    \includegraphics[width=0.8\linewidth]{\repodir/final/analysis/output/figures/53x_params_ionization_vsCFDK.png} 
    \caption{Comparison of MWT $\Delta AS$ and $\Delta PD$ versus CFD-predicted K for the 53x test case sequence. The color scale shows the nominal K mass fraction to indicate the order of the data points (the 53x sequences were an increasing 'ramp' in K mass fraction).}
    \label{fig:SI_53x_params_ionization_vsCFDK}
\end{figure}

\begin{figure}[h]
    \centering
    \includegraphics[width=0.8\linewidth]{\repodir/final/analysis/output/figures/53x_params_ionization_vsCFDKOH.png} 
    \caption{Comparison of MWT $\Delta AS$ and $\Delta PD$ versus CFD-predicted KOH for the 53x test case sequence. The color scale shows the nominal K mass fraction to indicate the order of the data points (the 53x sequences were an increasing 'ramp' in K mass fraction).}
    \label{fig:SI_53x_params_ionization_vsCFDKOH}
\end{figure}



\begin{figure}[]
\centering
\includegraphics[width=0.7\textwidth]{\repodir/experiment/analysis/mwt/output/figures/mwt_536_pos_mnum_sample.png}
\caption{MWT $AS$ time profiles for a sampling of individual acquisitions (indexed by 'measurement number': $mnum$) for different positions for the 536\_pos test case sequence.}
\label{fig:SI_mwt_536_pos_mnum_sample}
\end{figure}

\begin{figure}
    \centering
    \includegraphics[width=0.8\linewidth]{\repodir/experiment/analysis/mwt/output/figures/mwt_kwt_mnum_sample.png} 
    \caption{MWT $\Delta AS$ time profiles for individual acquisitions (indexed by 'measurement number': $mnum$) for different runs of the 53x test case.  }
    \label{fig:SI_mwt_kwt_mnum_sample}
\end{figure}

\begin{figure}
    \centering
    \includegraphics[width=0.8\linewidth]{\repodir/experiment/analysis/mwt/output/figures/mwt_53x_std_individual_count.png} 
    \caption{MWT $\Delta AS$ time profiles for individual test cases for the 53x test case. The shaded region represents the standard deviation across individual acquisitions, and the number of acquisitions is shown as a number in the plot.}
    \label{fig:SI_mwt_53x_std_individual_count}
\end{figure}

\begin{figure}
    \centering
    \includegraphics[]{\repodir/final/analysis/mwt/output/figures/AS_fft.png} 
    \caption{Power spectral density of the AS signal in the time window of 50 to 1 $\mu s$ before the laser pulse.}
    \label{fig:SI_AS_fft}
\end{figure}

\begin{figure}
    \centering
    \includegraphics[width=0.8\linewidth]{\repodir/experiment/analysis/auto/output/mwt_1d_noise/536_pos/hist2/2023-05-18_1.png} 
    \caption{Histogram of the magnitude of the MWT signal for different positions (indicated in mm at the top of the plot) for the 536\_pos test case for 2023-05-18 Run 1. The color represents the number of times the magnitude was observed within a magnitude bin (bin sizes are observable as pixels in the image). }
    \label{fig:SI_mwt_1d_noise_536_pos_hist2_2023-05-18_1}
\end{figure}

\begin{figure}
    \centering
    \includegraphics[width=0.8\linewidth]{\repodir/experiment/analysis/auto/output/mwt_1d_noise/536_pos/hist4/hist4_536_pos.png} 
    \caption{Histogram of the magnitude of the MWT signal for different positions for the 536\_pos test case sequence for each experiment Run. The y-axis is the number of times the magnitude was observed within a magnitude bin (same bins as in Figure \ref{fig:SI_mwt_1d_noise_536_pos_hist2_2023-05-18_1}). The x-axis is the magnitude of the MWT signal.}
    \label{fig:SI_mwt_1d_noise_536_pos_hist4_hist4_536_pos}
\end{figure}

\begin{figure}[]
\centering
\includegraphics[]{\repodir/final/analysis/mwt/output/figures/dAS_max_notorch_compare.png}
\caption{Comparison of $\Delta AS_{max}$ between multiple seed conditions: No torch, torch on and no seed (5x0), 1\% seed (5x6), and 0.1\% seed (5x3). All measurements were taken with an equivalence ratio of 0.8.}
\end{figure}


\begin{figure}
\centering
\includegraphics[]{\repodir/experiment/analysis/mwt/output/figures/mwt_fitting_compare.png}
\caption{Comparison of the fitting methods for the mono-molecular regime zoomed in to see the difference in time constant. }
\label{fig:SI_mwt_fitting_compare}
\end{figure}

\clearpage
\subsection{AES and MWT Seed and Position dependence fits}

This section contains plots of fits which led to the plotted fit parameters in the final text. In the following plots, the fitted parameters are averaged over all runs; however, during the fitting procedure, fits are performed on each individual run. 

\begin{figure}[]
\centering
\includegraphics[width=0.7\textwidth]{\repodir/final/dataset/output/figures/536_pos_phi08_absem_fit.png}
\caption{AES absorption $\alpha$ and fits for 536\_pos and 0.8 equivalence ratio. Note that some stage position values that were only performed on 2023-05-12 have been dropped to reduce the size of the figure. }
\label{fig:SI_536_pos_phi08_absem_fit}
\end{figure}

\begin{figure}[]
\centering
\includegraphics[width=0.7\textwidth]{\repodir/final/dataset/output/figures/516_pos_phi06_absem_fit.png}
\caption{$\alpha$ and fits for 536\_pos and 0.6 equivalence ratio. These measurements were only performed on 2023-05-24.}
\label{fig:SI_536_pos_phi06_absem_fit}
\end{figure}

\begin{figure}[]
\centering
\includegraphics[width=0.7\textwidth]{\repodir/final/dataset/output/figures/53x_absem_fit.png}
\caption{$\alpha$ and fits for 53x.}
\label{fig:SI_53x_absem_fit}
\end{figure}


\begin{figure}[]
\centering
\includegraphics[width=0.7\textwidth]{\repodir/final/dataset/output/figures/53x_mwt_fit_exp.png}
\caption{MWT measurements and fits for 53x for the exponential fitting method. }
\label{fig:SI_53x_mwt_fit_exp}
\end{figure}
